\begin{question}{15}{
    De Explosieven Opruimingsdienst Defensie (EOD) onderzoekt de vertragingstijd van ontstekingsmechanismen van onontplofte bommen uit de Tweede Wereldoorlog.
    Oorspronkelijk is deze vertragingstijd ontworpen dat een bom gemiddeld na $\mu_0 = 120$ milliseconden ontploft.
    Het risico bestaat dat door decennialange corrosie de gemiddelde vertragingstijd significant is veranderd.

    In een gecontroleerde testomgeving heeft de EOD een steekproef van $49$ WO2-bommen tot ontploffing gebracht, waarbij een gemiddelde van $117.5$ milliseconden en een standaardafwijking van $3.2$ milliseconden is waargenomen.
    Je mag hierbij aannemen dat de vertragingstijden normaal verdeeld zijn.
}

    \subquestion{8}{
        Bereken een \SI{98}{\percent}-betrouwbaarheidsinterval voor de gemiddelde vertragingstijd $\mu$ op basis van de steekproefgegevens.
        Rond de grenzen van het betrouwbaarheidsinterval af op \'e\'en decimaal, zodanig dat de betrouwbaarheid gewaarborgd blijft.
        Is er reden om aan te nemen dat de gemiddelde vertragingstijd significant is veranderd ten opzichte van de oorspronkelijke waarde van $\mu_0 = 120$ milliseconden?
    }

    \solution{
        In de vraag is het volgende gegeven: 
        \begin{itemize}
            \item de oorspronkelijke waarde $\mu_0 = 120$ milliseconden, 
            \item het steekproefgemiddelde $\samplemeanx = 117.5$,
            \item de steekproefstandaardafwijking $\samplestd = 3.2$, en
            \item de steekproefgrootte $n = 49$.
        \end{itemize}

        De vertragingstijd $X$ van een willekeurige WO2-bom is normaal verdeeld met $\mu = \text{ ?}$ en $\sigma = \text{ ?}$. 
        We willen een betrouwbaarheidsinterval voor de onbekende $\mu$ berekenen op basis van de steekproefgegevens. 
        
        Omdat de standaardafwijking $\sigma$ ook onbekend is, moeten we die schatten met $\samplestd = 3.2$ en voor het betrouwbaarheidsinterval moeten we gebruik maken van de $t$-verdeling.\rubric{1}

        Verder is $\alpha = 0.02$, omdat we een \SI{98}{\percent}-betrouwbaarheidsinterval willen berekenen.\rubric{1}

        De bijbehorende $t$-waarde is gelijk aan
        \[
            t = \invt(\opp=1-\frac{\alpha}{2}, \df=n-1) = \invt(\opp=1 - \frac{0.02}{2} = 0.99, \df=48) \approx 2.4066.\rubric{2}
        \]
        Het \SI{98}{\percent}-betrouwbaarheidsinterval wordt dan gegeven door
        \begin{align*}
            &\phantom{=}[\samplemeanx - t \cdot \frac{ \samplestd }{ \sqrt{ n } }, \samplemeanx + t \cdot \frac{ \samplestd }{ \sqrt{ n } }] \\
            &= [117.5 - 2.4066 \cdot \frac{ 3.2 }{ \sqrt{ 49 } }, 117.5 + 2.4066 \cdot \frac{ 3.2 }{ \sqrt{ 49 } }] \\
            &= [116.3998, 118.6002].\rubric{2}
        \end{align*}
        Om de betrouwbaarheid te waarborgen, moeten we het interval naar buiten afronden, oftewel op \'e\'en decimaal afgerond krijgen we een interval van $[116.3, 118.7]$ milliseconden. \rubric{1}

        Omdat de oorspronkelijke waarde $\mu_0 = 120$ milliseconden niet in het betrouwbaarheidsinterval ligt, is er voldoende reden om aan te nemen dat de gemiddelde vertragingstijd significant is veranderd ten opzichte van de oorspronkelijke waarde. \rubric{1}
    }

    \subquestion{7}{
        De EOD wil een nauwkeurigere schatting bepalen van de gemiddelde vertragingstijd van WO2-bommen, namelijk met een marge van $\pm 0.5$ milliseconden.
        Wat is de minimale steekproefomvang $n$ waarvoor het \SI{98}{\percent}-betrouwbaarheidsinterval voor $\mu$ maximaal deze afwijking heeft van het gemiddelde?
        Je mag aannemen dat de steekproefstandaardafwijking nog steeds gelijk is aan $3.2$ milliseconden.
    }
    \solution{
        Het betrouwbaarheidsinterval voor $\mu$ wordt gegeven door
        \[
            [\samplemeanx - t \cdot \samplestd / \sqrt{n}, \quad \samplemeanx + t \cdot \samplestd / \sqrt{n}],
        \]
        De afwijking van het betrouwbaarheidsinterval ten opzichte van het steekproefgemiddelde is dus $\pm\ t \cdot \samplestd / \sqrt{n}$ en we willen dat deze kleiner is dan of gelijk aan $\pm 0.5$ milliseconden.\rubric{2}

        Hiertoe voeren we op de grafische rekenmachine het volgende in in het functiescherm:
        \[
            y_1 = \frac{s}{\sqrt{X}} \cdot \invt(\opp=1-\alpha/2, \df=X-1) = \frac{3.2}{\sqrt{X}} \cdot \invt(\opp=0.99, \df=X-1) \rubric{2}
        \]
        Met behulp van de tabelfunctie krijgen we dan als volgt:
        \begin{align*}
            X = 224: &\quad \frac{3.2}{\sqrt{224}} \cdot \invt(\opp=0.99, \df=224-1) \approx 0.5010. \\
            X = 225: &\quad \frac{3.2}{\sqrt{225}} \cdot \invt(\opp=0.99, \df=225-1) \approx 0.4999. \rubric{2}
        \end{align*}
        Om een betrouwbaarheidsinterval te vinden met de gewenste precisie van maximaal $\pm\ 0.5$ meter, is een minimale steekproefgrootte van $n = 225$ WO2-bommen benodigd.\rubric{1}
    }
\end{question}
